\Page{3}{\textenglish{ABSTRACT}}{אופנה פוגשת בינה מלאכותית}{36.5}{43}

{\fontsize{13.1}{18.7}\selectfont

אופנה וטכנולוגיה נראות במבט ראשון כמו שני עולמות שמדברים בשפות שונות. אופנה מזוהה עם אמנות, יצירתיות ואסתטיקה, עם הרגש שמעורר צבע ועם הדרך שבה אנחנו מביעות את עצמנו. טכנולוגיה מזוהה עם מדע והיגיון, עם אלגוריתמים שמחפשים סדר ועם נתונים שאפשר להשוות. קל להבין למה נדמה שהמרחק ביניהן גדול. כשאנחנו מתאהבות במראה מסוים, המחשבה הראשונה היא בדרך כלל על מה שהוא מעורר בנו. קשה יותר לראות את שכבות העבודה הטכנולוגית שאפשרו למראה הזה להגיע לעינינו.

בתוך תעשיית האופנה, הקרבה בין העולמות האלה מתגלה כמעט בכל מקום. סטודיו משתמש בתוכנה כדי לתת צורה לרעיון שעדיין נמצא בראש של מעצבת. חזאית קוראת תצלומים מן הרחוב לצד נתוני חיפוש ומסחר. צוות קמפיין יכול להפיק תמונה בסביבה שקיימת רק על המסך. אתר קניות יכול להציג ללקוחה הדמיה של לוק על התמונה שלה. כל שימוש עונה על צורך אחר, והשינוי המצטבר רחב: הטכנולוגיה משפיעה על מה שמעצבות יכולות להציע, על מה שמותגים מעזים להזמין, על מה שמופיע בפרסומת ועל הדרכים שבהן אנחנו מגלות אופנה.

הכתבה על \textbf{עיצוב ופיתוח בגדים} נכנסת אל הרגע שבו רעיון מתחיל לקבל צורה שאפשר לבחון. אחרי היכרות קצרה עם הדרך המקובלת מסקיצה לדוגמה, נראה איך כלים דיגיטליים ובינה מלאכותית מרחיבים את מרחב הניסוי של הסטודיו. מעצבת יכולה לבחון כמה כיוונים כשהבחירה עדיין פתוחה, וצוות פיתוח יכול לנסח מוקדם יותר מה הוא מבקש מהיצרן. הערך היצירתי מתגלה באפשרות לתת הזדמנות גם לרעיון שהיה קשה להסביר או לבחון במהירות.

הכתבה על \textbf{חיזוי טרנדים} מתעניינת במקצוע שמנסה להבין מה הקהל ירצה ללבוש כשעוד מוקדם להזמין את העונה. נלמד לקרוא את הרמזים שחזאיות אוספות מתצוגות, מחנויות ומתרבות הרחוב, ונראה מה קורה כשאפשר לזהות דפוסים גם בתמונות ובמידע מהמסחר. מבחינת קניינית, ההבדל בין צבע שמעורר עניין לבין צבע שאנשים יקנו יכול לקבוע כמה להזמין ומתי להציע אותו. התוכנה מוסיפה דרכים לראות את ההבדל הזה. לחזאית יש עדיין אחריות לפרש אותו בשוק ובקהל מסוימים.

בכתבה על \textbf{קמפיינים ודמויות וירטואליות} ניכנס אל העבודה שמאחורי תמונה שנראית כאילו נולדה ברגע אחד. יום צילום נשען על ליהוק, סטיילינג, איפור, תאורה, מקום ותיאום בין אנשים. בהפקה שמשלבת בינה מלאכותית, חלק מהעבודה עובר להכנת חומרי ייחוס ולבחירה בין גרסאות. לצד דוגמניות שמופיעות באמצעות ייצוג דיגיטלי שלהן, יש גם דמויות שנבנו כדי להמשיך להופיע בפיד במשך זמן. ההבדל משנה את העבודה הדרושה לצוות ואת מה שהקהל מצפה לקבל מהתמונה או מהחשבון.

הכתבה על \textbf{עתיד האופנה} תתחיל במה שכבר אפשר לנסות ותפנה אל האפשרויות שעוד מתגבשות. הדמיית לבוש וקנייה מודרכת עשויות לשנות את הדרך שבה אנחנו מגלות פריטים. כלים לתכנון ביקוש עשויים לעזור למותגים להחליט מה לייצר ובאיזו כמות, והשאלה הסביבתית תיבחן לפי מה שקורה לעודפים בפועל. גם העבודה עצמה עשויה להשתנות, משילוב של פיתוח דיגיטלי עם ידע בחומרים ועד לתפקידים חדשים סביב הפקת דימויים ובדיקת התאמה. יש כאן עתיד שאפשר להתחיל לדמיין מתוך חוויות שכבר מוכרות לנו: להיכנס לחנות, לגלות מראה חדש ולבחור מה ראוי למקום בארון.

\par}
